\section{Related Work}\label{sec:related_work}

\subsection{VLA Robustness}

Methods for improving VLA robustness divide into training-time and inference-time approaches.

\paragraph{Training-time methods.}
RobustVLA~\cite{robustvla2025} combines RL post-training with Jacobian regularization and smoothness regularization to stabilize action predictions under observation noise and action perturbations.
SRPO~\cite{srpo2025} applies self-referential RL fine-tuning, using the policy's own successful rollouts as references to shape rewards for failed trajectories via latent world model representations.
DiG-Flow~\cite{digflow2025} uses distributional discrepancy between observation and action embeddings as a gating signal to rebalance representation-level control, without modifying the flow-matching trajectory.
TT-VLA~\cite{liu2026ttvla} blurs the training/inference boundary by adapting VLA parameters at test time through RL with progress rewards, but still requires per-task reward design and gradient updates during deployment.
A shared limitation of these methods is the requirement for retraining, reward engineering, or per-task RL---making them difficult to apply as drop-in improvements to an arbitrary frozen VLA.

\paragraph{Inference-time methods.}
ReconVLA~\cite{chen2026reconvla} applies conformal quantile regression to calibrate per-action uncertainty bounds and uses squared Mahalanobis distance for out-of-distribution detection, enabling uncertainty-guided action selection and proactive failure detection without retraining.
SC-VLA~\cite{liu2026scvla} performs online action refinement through sparse world imagination combined with residual RL, but requires task-level RL training for the correction policy.
VLA-Pilot~\cite{li2025vlapilot} applies evolutionary diffusion optimization guided by an MLLM verifier, achieving plug-and-play correction at the cost of 2.41\,s per control step---prohibitive for real-time manipulation.
A2C2~\cite{a2c2_2025} adds a lightweight asynchronous correction head to a frozen VLA that runs in parallel with chunk generation, correcting each action within a pre-generated chunk to restore closed-loop responsiveness; however, it applies correction unconditionally without uncertainty-guided selection or explicit magnitude control.

A2C2 is the closest concurrent work to ACR-VLA: both attach a lightweight correction module to a frozen VLA without retraining the base policy.
The key architectural difference is that ACR-VLA operates as a two-stage pipeline---multi-sample conformal selection followed by bounded residual correction---whereas A2C2 applies per-step correction directly to the base policy's action chunks without a selection stage.
ACR-VLA additionally enforces explicit magnitude control via $L_2$ norm clipping, which our ablations show is critical for preventing compounding correction errors (Section~\ref{sec:magnitude_results}).
Selection alone is upper-bounded by sample quality---it cannot produce an action better than the best sample.
Correction alone lacks a mechanism to identify a good starting point.
ACR-VLA addresses both by selecting a low-nonconformity candidate via Mahalanobis scoring, then applying a bounded residual correction trained offline with supervised $L_2$ loss.

\subsection{Conformal Prediction in Robotics}

Conformal prediction (CP) provides distribution-free coverage guarantees under exchangeability~\cite{angelopoulos2023gentle}.
For multivariate outputs, Braun et al.~\cite{braun2025cgs} develop Conformalized Gaussian Scoring (CGS), which uses Mahalanobis-based nonconformity scores to construct prediction regions that respect cross-dimensional covariance---the scoring method ACR-VLA builds upon for its selection stage.

In robotics, CP has been applied to several safety-critical settings.
Conformalized Teleoperation~\cite{zhao2024conformalized} uses CP to quantify uncertainty in high-DOF action mappings during shared autonomy, flagging unreliable predictions for human review.
CP-in-the-Loop~\cite{wang2025cploop} integrates conformal prediction regions as feedback constraints in trajectory optimization, tightening or relaxing plans based on calibrated uncertainty.
ACP+CBF~\cite{zhou2025acp_cbf} combines adaptive conformal prediction with control barrier functions to maintain safety constraints under distribution shift in RL policies.
A common pattern across these methods is that CP serves as a \emph{binary decision mechanism}---accept or reject, safe or unsafe, execute or abstain.

ACR-VLA uses Mahalanobis nonconformity scores derived from shrinkage covariance estimates in its selection stage, but does \emph{not} claim formal coverage guarantees for the corrected output.
CP serves here as a scoring mechanism: the Mahalanobis score ranks samples for selection, and the covariance structure is available as conditioning features for the correction module (though our ablations find that observation features alone suffice; Section~\ref{sec:ablation}).
The downstream correction module produces a refined action whose relationship to the original conformal region is no longer governed by coverage theory.

\subsection{Residual Policy Learning}

Residual policy learning augments a base controller with a learned additive correction.
Silver et al.~\cite{silver2018residual} introduced the framework for classical controllers, using zero-initialization to ensure the learned residual starts from the base policy's behavior and gradually improves upon it.
Subsequent work has focused on controlling residual magnitude---a central design concern, since unconstrained corrections lead to compounding error and policy collapse.
TRANSIC~\cite{jiang2024transic} uses a learned binary gate conditioned on the current state to decide whether to apply a sim-to-real residual, demonstrating that ungated correction causes model collapse.
PLD~\cite{xiao2026pld} freezes the VLA backbone and trains lightweight residual actors with a magnitude bound $\xi$ tuned per benchmark via off-policy RL, then distills improved trajectories back into the base model.
DRC~\cite{pan2025drc} applies exponential decay to each expert correction event, allowing corrective signals to persist and fade after application.

Magnitude control is the key connection to ACR-VLA: Silver et al.\ use zero-initialization, TRANSIC uses binary gating, PLD uses per-benchmark magnitude bounds, DRC uses exponential decay after correction events.
ACR-VLA uses fixed $L_2$ norm clipping---the simplest possible scheme---yet our ablations show that the choice of clipping threshold $\varepsilon$ has an outsized effect on performance, with the optimal value varying across task suites (Section~\ref{sec:experiments}).
We view this as evidence that magnitude control in VLA inference-time correction is an underexplored hyperparameter.
A second distinction separates ACR-VLA from the above methods: all four require RL training (reward functions, environment interaction, or on-policy rollouts), whereas ACR-VLA's correction module is trained with supervised $L_2$ loss on offline demonstration data, requiring neither reward design nor environment access during training.

\subsection{Uncertainty Quantification for VLAs}

Recent work has developed UQ methods tailored to VLA architectures.
DiU~\cite{xing2026diu} derives a closed-form posterior covariance for flow-matching models by computing the divergence of the velocity field, enabling single-pass per-DOF uncertainty estimation without additional sampling.
SAFE~\cite{gu2025safe} trains a multitask failure detector on VLA internal features and uses conformal prediction to calibrate detection timing, triggering intervention when anomalous states are detected.
Zollo and Zemel~\cite{zollo2025vla_calibration} provide the first systematic study of confidence calibration in VLAs, finding that prompt ensembles combined with Platt scaling improve calibration across manipulation tasks.

ACR-VLA takes a multi-sample approach to UQ ($K{=}10$ forward passes), in contrast to DiU's single-pass estimation.
The multi-sample strategy is more expensive but captures both epistemic and aleatoric uncertainty components---a property that matters under distribution shift, where epistemic uncertainty rises but single-pass aleatoric estimates may not.
Functionally, ACR-VLA uses uncertainty for action selection (identifying the lowest-nonconformity sample), whereas DiU targets efficient uncertainty visualization, SAFE targets failure detection, and VLA calibration~\cite{zollo2025vla_calibration} targets confidence reliability.
Our ablations reveal a negative result connecting to this literature: per-DOF conformal features, despite their richer structure, do not improve correction performance when used as CRM input at the $K{=}10$ budget (Section~\ref{sec:ablation}), suggesting that translating fine-grained uncertainty estimates into actionable correction signals remains an open challenge.
